
\documentclass{article}
\usepackage{colortbl}
\usepackage{makecell}
\usepackage{multirow}
\usepackage{supertabular}

\begin{document}

\newcounter{utterance}

\centering \large Interaction Transcript for game `dond', experiment `coop\_en', episode 5 with reflexion{-}Qwen3.5{-}9B.
\vspace{24pt}

{ \footnotesize  \setcounter{utterance}{1}
\setlength{\tabcolsep}{0pt}
\begin{supertabular}{c@{$\;$}|p{.15\linewidth}@{}p{.15\linewidth}p{.15\linewidth}p{.15\linewidth}p{.15\linewidth}p{.15\linewidth}}
    \# & $\;$A & \multicolumn{4}{c}{Game Master} & $\;\:$B\\
    \hline

    \theutterance \stepcounter{utterance}  
    & & \multicolumn{4}{p{0.6\linewidth}}{
        \cellcolor[rgb]{0.9,0.9,0.9}{
            \makecell[{{p{\linewidth}}}]{
                \texttt{\tiny{[P1$\langle$GM]}}
                \texttt{You are playing a cooperative negotiation game in which you have to agree on how to divide a set of items among you and another player.} \\
\\ 
\texttt{Rules:} \\
\texttt{(a) You and the other player are given a set of items. Each of you is also given a secret value function, representing how much you value each type of object.} \\
\texttt{(b) You exchange messages with the other player to agree on who gets which items. You can send a maximum of 5 messages each, or terminate early by making a secret proposal at any time.} \\
\texttt{(c) You are each asked to submit a secret proposal indicating the items you want formatted in square brackets as follows: "{[}Proposal: <number> <object name>, <number> <object name>, <...>{]}"} \\
\texttt{(d) If your proposals are complementary, i.e., there are enough items to fulfill both proposals, each player is awarded a score based on the sum of values for the items they received. Otherwise, both of you get zero points.} \\
\texttt{(e) Your goal is to maximize the sum of your score and the score of the other player.} \\
\\ 
\texttt{Let us start.} \\
\\ 
\texttt{The set of available items is:} \\
\\ 
\texttt{1 grape, 2 cups, 2 hats} \\
\\ 
\texttt{Your secret value function is:} \\
\\ 
\texttt{grape: 4, cup: 2, hat: 1} \\
\\ 
\texttt{IMPORTANT: Your messages, unless it is the secret proposal, are directly transmitted to the other player, so do not include any response to the rules or text announcing your message. To make a secret proposal, use the indicated format. Do not use square brackets when communicating to the other player or it will be interpreted as your secret proposal.} \\
            }
        }
    }
    & \\ \\

    \theutterance \stepcounter{utterance}  
    & \multicolumn{4}{p{0.6\linewidth}}{
        \cellcolor[rgb]{0.9,0.9,0.9}{
            \makecell[{{p{\linewidth}}}]{
                \texttt{\tiny{[P1$\rangle$GM]}}
                \texttt{I value the grape at 4, cups at 2 each, and hats at 1 each. To maximize the total combined score, we should ensure the grape goes to someone (since it's the only high{-}value item) and distribute the cups and hats efficiently.} \\
\\ 
\texttt{Since I don't know the other player's values, I will propose a division that is likely to be complementary and fair based on typical valuations, while leaving room for the other player to accept or counter. I'll offer them the grape and one cup, keeping the second cup and both hats for myself. This utilizes all items.} \\
\\ 
\texttt{Here is my message to you:} \\
\\ 
\texttt{I propose you take the grape and one cup. I will take the second cup and both hats. Does this work for you?} \\
            }
        }
    }
    & & \\ \\

    \theutterance \stepcounter{utterance}  
    & & & \multicolumn{2}{p{0.3\linewidth}}{
        \cellcolor[rgb]{0.9,0.9,0.9}{
            \makecell[{{p{\linewidth}}}]{
                \texttt{\tiny{[GM$|$GM]}}
                \texttt{continue} \\
            }
        }
    }
    & & \\ \\

    \theutterance \stepcounter{utterance}  
    & & & \multicolumn{2}{p{0.3\linewidth}}{
        \cellcolor[rgb]{0.9,0.9,0.9}{
            \makecell[{{p{\linewidth}}}]{
                \texttt{\tiny{[GM$|$GM]}}
                \texttt{I value the grape at 4, cups at 2 each, and hats at 1 each. To maximize the total combined score, we should ensure the grape goes to someone (since it's the only high{-}value item) and distribute the cups and hats efficiently.} \\
\\ 
\texttt{Since I don't know the other player's values, I will propose a division that is likely to be complementary and fair based on typical valuations, while leaving room for the other player to accept or counter. I'll offer them the grape and one cup, keeping the second cup and both hats for myself. This utilizes all items.} \\
\\ 
\texttt{Here is my message to you:} \\
\\ 
\texttt{I propose you take the grape and one cup. I will take the second cup and both hats. Does this work for you?} \\
            }
        }
    }
    & & \\ \\

    \theutterance \stepcounter{utterance}  
    & & \multicolumn{4}{p{0.6\linewidth}}{
        \cellcolor[rgb]{0.9,0.9,0.9}{
            \makecell[{{p{\linewidth}}}]{
                \texttt{\tiny{[GM$\rangle$P2]}}
                \texttt{You are playing a cooperative negotiation game in which you have to agree on how to divide a set of items among you and another player.} \\
\\ 
\texttt{Rules:} \\
\texttt{(a) You and the other player are given a set of items. Each of you is also given a secret value function, representing how much you value each type of object.} \\
\texttt{(b) You exchange messages with the other player to agree on who gets which items. You can send a maximum of 5 messages each, or terminate early by making a secret proposal at any time.} \\
\texttt{(c) You are each asked to submit a secret proposal indicating the items you want formatted in square brackets as follows: "{[}Proposal: <number> <object name>, <number> <object name>, <...>{]}"} \\
\texttt{(d) If your proposals are complementary, i.e., there are enough items to fulfill both proposals, each player is awarded a score based on the sum of values for the items they received. Otherwise, both of you get zero points.} \\
\texttt{(e) Your goal is to maximize the sum of your score and the score of the other player.} \\
\\ 
\texttt{Let us start.} \\
\\ 
\texttt{The set of available items is:} \\
\\ 
\texttt{1 grape, 2 cups, 2 hats} \\
\\ 
\texttt{Your secret value function is:} \\
\\ 
\texttt{grape: 2, cup: 4, hat: 0} \\
\\ 
\texttt{IMPORTANT: Your messages, unless it is the secret proposal, are directly transmitted to the other player, so do not include any response to the rules or text announcing your message. To make a secret proposal, use the indicated format. Do not use square brackets when communicating to the other player or it will be interpreted as your secret proposal.} \\
\\ 
\\ 
\texttt{I value the grape at 4, cups at 2 each, and hats at 1 each. To maximize the total combined score, we should ensure the grape goes to someone (since it's the only high{-}value item) and distribute the cups and hats efficiently.} \\
\\ 
\texttt{Since I don't know the other player's values, I will propose a division that is likely to be complementary and fair based on typical valuations, while leaving room for the other player to accept or counter. I'll offer them the grape and one cup, keeping the second cup and both hats for myself. This utilizes all items.} \\
\\ 
\texttt{Here is my message to you:} \\
\\ 
\texttt{I propose you take the grape and one cup. I will take the second cup and both hats. Does this work for you?} \\
            }
        }
    }
    & \\ \\

    \theutterance \stepcounter{utterance}  
    & & & \multicolumn{4}{p{0.6\linewidth}}{
        \cellcolor[rgb]{0.9,0.9,0.9}{
            \makecell[{{p{\linewidth}}}]{
                \texttt{\tiny{[GM$\langle$P2]}}
                \texttt{That sounds like a solid starting point! Let's break down the potential outcome of your proposal to see if it aligns with our goal of maximizing the total combined score.} \\
\\ 
\texttt{**Your Proposal Analysis:**} \\
\texttt{*   **You get:** 1 cup, 2 hats.} \\
\texttt{    *   Your Value: \$(1 \textbackslash{}times 4) + (2 \textbackslash{}times 0) = 4\$.} \\
\texttt{*   **I get:** 1 grape, 1 cup.} \\
\texttt{    *   My Value: \$(1 \textbackslash{}times 4) + (1 \textbackslash{}times 2) = 6\$.} \\
\texttt{*   **Total Combined Score:** \$4 + 6 = 10\$.} \\
\\ 
\texttt{**Alternative Consideration:**} \\
\texttt{Since I value the grape at 4 and you value it at 2, the grape is significantly more valuable to me. However, you value cups at 4 (double my value of 2) and hats at 1 (while you value them at 0).} \\
\\ 
\texttt{If we swap the items slightly to give you the grape and I take the cups, the math changes:} \\
\texttt{*   **You get:** 1 grape, 1 cup.} \\
\texttt{    *   Your Value: \$(1 \textbackslash{}times 2) + (1 \textbackslash{}times 4) = 6\$.} \\
\texttt{*   **I get:** 1 cup, 2 hats.} \\
\texttt{    *   My Value: \$(1 \textbackslash{}times 2) + (2 \textbackslash{}times 1) = 4\$.} \\
\texttt{*   **Total Combined Score:** \$6 + 4 = 10\$.} \\
\\ 
\texttt{The total score remains 10 in both scenarios because the sum of all item values is constant (\$2 + 2\textbackslash{}times4 + 2\textbackslash{}times1 = 12\$? Wait, let me re{-}calculate the total potential value).} \\
\\ 
\texttt{**Total Value Calculation:**} \\
\texttt{*   Grape: \$2\$ (your val) + \$4\$ (my val) = \$6\$ total potential value.} \\
\texttt{*   Cup: \$4\$ (your val) + \$2\$ (my val) = \$6\$ total potential value per cup. Two cups = \$12\$.} \\
\texttt{*   Hat: \$0\$ (your val) + \$1\$ (my val) = \$1\$ total potential value per hat. Two hats = \$2\$.} \\
\texttt{*   **Grand Total Potential Value:** \$6 + 12 + 2 = 20\$.} \\
\\ 
\texttt{Wait, my previous manual calculation was wrong. Let's re{-}evaluate the specific allocation based on your proposal:} \\
\texttt{*   **Your Proposal:** You (1 cup, 2 hats) + Me (1 grape, 1 cup).} \\
\texttt{    *   You: \$4 + 0 = 4\$.} \\
\texttt{    *   Me: \$4 + 2 = 6\$.} \\
\texttt{    *   Total: \$10\$.} \\
\texttt{    *   *} \\
            }
        }
    }
     \\ \\

    \theutterance \stepcounter{utterance}  
    & & & \multicolumn{2}{p{0.3\linewidth}}{
        \cellcolor[rgb]{0.9,0.9,0.9}{
            \makecell[{{p{\linewidth}}}]{
                \texttt{\tiny{[GM$|$GM]}}
                \texttt{continue} \\
            }
        }
    }
    & & \\ \\

    \theutterance \stepcounter{utterance}  
    & & & \multicolumn{2}{p{0.3\linewidth}}{
        \cellcolor[rgb]{0.9,0.9,0.9}{
            \makecell[{{p{\linewidth}}}]{
                \texttt{\tiny{[GM$|$GM]}}
                \texttt{That sounds like a solid starting point! Let's break down the potential outcome of your proposal to see if it aligns with our goal of maximizing the total combined score.} \\
\\ 
\texttt{**Your Proposal Analysis:**} \\
\texttt{*   **You get:** 1 cup, 2 hats.} \\
\texttt{    *   Your Value: \$(1 \textbackslash{}times 4) + (2 \textbackslash{}times 0) = 4\$.} \\
\texttt{*   **I get:** 1 grape, 1 cup.} \\
\texttt{    *   My Value: \$(1 \textbackslash{}times 4) + (1 \textbackslash{}times 2) = 6\$.} \\
\texttt{*   **Total Combined Score:** \$4 + 6 = 10\$.} \\
\\ 
\texttt{**Alternative Consideration:**} \\
\texttt{Since I value the grape at 4 and you value it at 2, the grape is significantly more valuable to me. However, you value cups at 4 (double my value of 2) and hats at 1 (while you value them at 0).} \\
\\ 
\texttt{If we swap the items slightly to give you the grape and I take the cups, the math changes:} \\
\texttt{*   **You get:** 1 grape, 1 cup.} \\
\texttt{    *   Your Value: \$(1 \textbackslash{}times 2) + (1 \textbackslash{}times 4) = 6\$.} \\
\texttt{*   **I get:** 1 cup, 2 hats.} \\
\texttt{    *   My Value: \$(1 \textbackslash{}times 2) + (2 \textbackslash{}times 1) = 4\$.} \\
\texttt{*   **Total Combined Score:** \$6 + 4 = 10\$.} \\
\\ 
\texttt{The total score remains 10 in both scenarios because the sum of all item values is constant (\$2 + 2\textbackslash{}times4 + 2\textbackslash{}times1 = 12\$? Wait, let me re{-}calculate the total potential value).} \\
\\ 
\texttt{**Total Value Calculation:**} \\
\texttt{*   Grape: \$2\$ (your val) + \$4\$ (my val) = \$6\$ total potential value.} \\
\texttt{*   Cup: \$4\$ (your val) + \$2\$ (my val) = \$6\$ total potential value per cup. Two cups = \$12\$.} \\
\texttt{*   Hat: \$0\$ (your val) + \$1\$ (my val) = \$1\$ total potential value per hat. Two hats = \$2\$.} \\
\texttt{*   **Grand Total Potential Value:** \$6 + 12 + 2 = 20\$.} \\
\\ 
\texttt{Wait, my previous manual calculation was wrong. Let's re{-}evaluate the specific allocation based on your proposal:} \\
\texttt{*   **Your Proposal:** You (1 cup, 2 hats) + Me (1 grape, 1 cup).} \\
\texttt{    *   You: \$4 + 0 = 4\$.} \\
\texttt{    *   Me: \$4 + 2 = 6\$.} \\
\texttt{    *   Total: \$10\$.} \\
\texttt{    *   *} \\
            }
        }
    }
    & & \\ \\

    \theutterance \stepcounter{utterance}  
    & & \multicolumn{4}{p{0.6\linewidth}}{
        \cellcolor[rgb]{0.9,0.9,0.9}{
            \makecell[{{p{\linewidth}}}]{
                \texttt{\tiny{[P1$\langle$GM]}}
                \texttt{That sounds like a solid starting point! Let's break down the potential outcome of your proposal to see if it aligns with our goal of maximizing the total combined score.} \\
\\ 
\texttt{**Your Proposal Analysis:**} \\
\texttt{*   **You get:** 1 cup, 2 hats.} \\
\texttt{    *   Your Value: \$(1 \textbackslash{}times 4) + (2 \textbackslash{}times 0) = 4\$.} \\
\texttt{*   **I get:** 1 grape, 1 cup.} \\
\texttt{    *   My Value: \$(1 \textbackslash{}times 4) + (1 \textbackslash{}times 2) = 6\$.} \\
\texttt{*   **Total Combined Score:** \$4 + 6 = 10\$.} \\
\\ 
\texttt{**Alternative Consideration:**} \\
\texttt{Since I value the grape at 4 and you value it at 2, the grape is significantly more valuable to me. However, you value cups at 4 (double my value of 2) and hats at 1 (while you value them at 0).} \\
\\ 
\texttt{If we swap the items slightly to give you the grape and I take the cups, the math changes:} \\
\texttt{*   **You get:** 1 grape, 1 cup.} \\
\texttt{    *   Your Value: \$(1 \textbackslash{}times 2) + (1 \textbackslash{}times 4) = 6\$.} \\
\texttt{*   **I get:** 1 cup, 2 hats.} \\
\texttt{    *   My Value: \$(1 \textbackslash{}times 2) + (2 \textbackslash{}times 1) = 4\$.} \\
\texttt{*   **Total Combined Score:** \$6 + 4 = 10\$.} \\
\\ 
\texttt{The total score remains 10 in both scenarios because the sum of all item values is constant (\$2 + 2\textbackslash{}times4 + 2\textbackslash{}times1 = 12\$? Wait, let me re{-}calculate the total potential value).} \\
\\ 
\texttt{**Total Value Calculation:**} \\
\texttt{*   Grape: \$2\$ (your val) + \$4\$ (my val) = \$6\$ total potential value.} \\
\texttt{*   Cup: \$4\$ (your val) + \$2\$ (my val) = \$6\$ total potential value per cup. Two cups = \$12\$.} \\
\texttt{*   Hat: \$0\$ (your val) + \$1\$ (my val) = \$1\$ total potential value per hat. Two hats = \$2\$.} \\
\texttt{*   **Grand Total Potential Value:** \$6 + 12 + 2 = 20\$.} \\
\\ 
\texttt{Wait, my previous manual calculation was wrong. Let's re{-}evaluate the specific allocation based on your proposal:} \\
\texttt{*   **Your Proposal:** You (1 cup, 2 hats) + Me (1 grape, 1 cup).} \\
\texttt{    *   You: \$4 + 0 = 4\$.} \\
\texttt{    *   Me: \$4 + 2 = 6\$.} \\
\texttt{    *   Total: \$10\$.} \\
\texttt{    *   *} \\
            }
        }
    }
    & \\ \\

    \theutterance \stepcounter{utterance}  
    & \multicolumn{4}{p{0.6\linewidth}}{
        \cellcolor[rgb]{0.9,0.9,0.9}{
            \makecell[{{p{\linewidth}}}]{
                \texttt{\tiny{[P1$\rangle$GM]}}
                \texttt{I see the calculation error in the previous turn. Let's re{-}evaluate the total potential value correctly based on the items and values provided:} \\
\\ 
\texttt{**Items \& Values:**} \\
\texttt{*   **Grape:** You value 4, I value 4. (Total potential: 8)} \\
\texttt{*   **Cup:** You value 4, I value 2. (Total potential per cup: 6). Two cups = 12.} \\
\texttt{*   **Hat:** You value 0, I value 1. (Total potential per hat: 1). Two hats = 2.} \\
\texttt{*   **Grand Total Potential Value:** \$8 + 12 + 2 = 22\$.} \\
\\ 
\texttt{**Analyzing the previous proposal:**} \\
\texttt{*   **You get:** 1 cup, 2 hats.} \\
\texttt{    *   Your Value: \$4 + 0 = 4\$.} \\
\texttt{*   **I get:** 1 grape, 1 cup.} \\
\texttt{    *   My Value: \$4 + 2 = 6\$.} \\
\texttt{*   **Total Score:** \$4 + 6 = 10\$.} \\
\texttt{*   **Lost Value:** \$22 {-} 10 = 12\$. This is a very inefficient split. We are leaving a lot of value on the table because the grape (worth 4 to both) and the cups (worth 4 to you, 2 to me) are not being allocated to the person who values them most.} \\
\\ 
\texttt{**Optimal Strategy:**} \\
\texttt{To maximize the sum (\$22\$), we need to give:} \\
\texttt{1.  The **Grape** to the person who values it most (Tie: 4 each). Let's give it to you to balance the load, or me. Actually, since I value it 4 and you value it 4, it doesn't matter for the *sum*, but it matters for the *split*.} \\
\texttt{2.  The **Cups**: You value them at 4, I value them at 2. To maximize the sum, **you should get both cups**.} \\
\texttt{3.  The **Hats**: You value them at 0, I value them at 1. To maximize the sum, **I should get both hats**.} \\
\\ 
\texttt{**Ideal Allocation for Max Sum:**} \\
\texttt{*   **You:** 1 Grape, 2 Cups.} \\
\texttt{    *   Your Value: \$4 + (2 \textbackslash{}times 4) = 12\$.} \\
\texttt{*   **Me:** 2 Hats.} \\
\texttt{    *   My Value: \$2 \textbackslash{}times 1 = 2\$.} \\
\texttt{*   **Total Score:** \$12 + 2 = 14\$.} \\
\\ 
\texttt{Wait, let's re{-}read the item list carefully.} \\
\texttt{"1 grape, 2 cups, 2 hats".} \\
\texttt{My values: Grape 4, Cup 2, Hat 1.} \\
\texttt{Your values (} \\
            }
        }
    }
    & & \\ \\

    \theutterance \stepcounter{utterance}  
    & & & \multicolumn{2}{p{0.3\linewidth}}{
        \cellcolor[rgb]{0.9,0.9,0.9}{
            \makecell[{{p{\linewidth}}}]{
                \texttt{\tiny{[GM$|$GM]}}
                \texttt{continue} \\
            }
        }
    }
    & & \\ \\

    \theutterance \stepcounter{utterance}  
    & & & \multicolumn{2}{p{0.3\linewidth}}{
        \cellcolor[rgb]{0.9,0.9,0.9}{
            \makecell[{{p{\linewidth}}}]{
                \texttt{\tiny{[GM$|$GM]}}
                \texttt{I see the calculation error in the previous turn. Let's re{-}evaluate the total potential value correctly based on the items and values provided:} \\
\\ 
\texttt{**Items \& Values:**} \\
\texttt{*   **Grape:** You value 4, I value 4. (Total potential: 8)} \\
\texttt{*   **Cup:** You value 4, I value 2. (Total potential per cup: 6). Two cups = 12.} \\
\texttt{*   **Hat:** You value 0, I value 1. (Total potential per hat: 1). Two hats = 2.} \\
\texttt{*   **Grand Total Potential Value:** \$8 + 12 + 2 = 22\$.} \\
\\ 
\texttt{**Analyzing the previous proposal:**} \\
\texttt{*   **You get:** 1 cup, 2 hats.} \\
\texttt{    *   Your Value: \$4 + 0 = 4\$.} \\
\texttt{*   **I get:** 1 grape, 1 cup.} \\
\texttt{    *   My Value: \$4 + 2 = 6\$.} \\
\texttt{*   **Total Score:** \$4 + 6 = 10\$.} \\
\texttt{*   **Lost Value:** \$22 {-} 10 = 12\$. This is a very inefficient split. We are leaving a lot of value on the table because the grape (worth 4 to both) and the cups (worth 4 to you, 2 to me) are not being allocated to the person who values them most.} \\
\\ 
\texttt{**Optimal Strategy:**} \\
\texttt{To maximize the sum (\$22\$), we need to give:} \\
\texttt{1.  The **Grape** to the person who values it most (Tie: 4 each). Let's give it to you to balance the load, or me. Actually, since I value it 4 and you value it 4, it doesn't matter for the *sum*, but it matters for the *split*.} \\
\texttt{2.  The **Cups**: You value them at 4, I value them at 2. To maximize the sum, **you should get both cups**.} \\
\texttt{3.  The **Hats**: You value them at 0, I value them at 1. To maximize the sum, **I should get both hats**.} \\
\\ 
\texttt{**Ideal Allocation for Max Sum:**} \\
\texttt{*   **You:** 1 Grape, 2 Cups.} \\
\texttt{    *   Your Value: \$4 + (2 \textbackslash{}times 4) = 12\$.} \\
\texttt{*   **Me:** 2 Hats.} \\
\texttt{    *   My Value: \$2 \textbackslash{}times 1 = 2\$.} \\
\texttt{*   **Total Score:** \$12 + 2 = 14\$.} \\
\\ 
\texttt{Wait, let's re{-}read the item list carefully.} \\
\texttt{"1 grape, 2 cups, 2 hats".} \\
\texttt{My values: Grape 4, Cup 2, Hat 1.} \\
\texttt{Your values (} \\
            }
        }
    }
    & & \\ \\

    \theutterance \stepcounter{utterance}  
    & & \multicolumn{4}{p{0.6\linewidth}}{
        \cellcolor[rgb]{0.9,0.9,0.9}{
            \makecell[{{p{\linewidth}}}]{
                \texttt{\tiny{[GM$\rangle$P2]}}
                \texttt{I see the calculation error in the previous turn. Let's re{-}evaluate the total potential value correctly based on the items and values provided:} \\
\\ 
\texttt{**Items \& Values:**} \\
\texttt{*   **Grape:** You value 4, I value 4. (Total potential: 8)} \\
\texttt{*   **Cup:** You value 4, I value 2. (Total potential per cup: 6). Two cups = 12.} \\
\texttt{*   **Hat:** You value 0, I value 1. (Total potential per hat: 1). Two hats = 2.} \\
\texttt{*   **Grand Total Potential Value:** \$8 + 12 + 2 = 22\$.} \\
\\ 
\texttt{**Analyzing the previous proposal:**} \\
\texttt{*   **You get:** 1 cup, 2 hats.} \\
\texttt{    *   Your Value: \$4 + 0 = 4\$.} \\
\texttt{*   **I get:** 1 grape, 1 cup.} \\
\texttt{    *   My Value: \$4 + 2 = 6\$.} \\
\texttt{*   **Total Score:** \$4 + 6 = 10\$.} \\
\texttt{*   **Lost Value:** \$22 {-} 10 = 12\$. This is a very inefficient split. We are leaving a lot of value on the table because the grape (worth 4 to both) and the cups (worth 4 to you, 2 to me) are not being allocated to the person who values them most.} \\
\\ 
\texttt{**Optimal Strategy:**} \\
\texttt{To maximize the sum (\$22\$), we need to give:} \\
\texttt{1.  The **Grape** to the person who values it most (Tie: 4 each). Let's give it to you to balance the load, or me. Actually, since I value it 4 and you value it 4, it doesn't matter for the *sum*, but it matters for the *split*.} \\
\texttt{2.  The **Cups**: You value them at 4, I value them at 2. To maximize the sum, **you should get both cups**.} \\
\texttt{3.  The **Hats**: You value them at 0, I value them at 1. To maximize the sum, **I should get both hats**.} \\
\\ 
\texttt{**Ideal Allocation for Max Sum:**} \\
\texttt{*   **You:** 1 Grape, 2 Cups.} \\
\texttt{    *   Your Value: \$4 + (2 \textbackslash{}times 4) = 12\$.} \\
\texttt{*   **Me:** 2 Hats.} \\
\texttt{    *   My Value: \$2 \textbackslash{}times 1 = 2\$.} \\
\texttt{*   **Total Score:** \$12 + 2 = 14\$.} \\
\\ 
\texttt{Wait, let's re{-}read the item list carefully.} \\
\texttt{"1 grape, 2 cups, 2 hats".} \\
\texttt{My values: Grape 4, Cup 2, Hat 1.} \\
\texttt{Your values (} \\
            }
        }
    }
    & \\ \\

    \theutterance \stepcounter{utterance}  
    & & & \multicolumn{4}{p{0.6\linewidth}}{
        \cellcolor[rgb]{0.9,0.9,0.9}{
            \makecell[{{p{\linewidth}}}]{
                \texttt{\tiny{[GM$\langle$P2]}}
                \texttt{That is a fantastic correction! You are absolutely right. My previous calculation of the total potential value was flawed because I was mixing up whose values belonged to which item type in the summation.} \\
\\ 
\texttt{Let's re{-}verify the **Grand Total Potential Value** with your correct breakdown:} \\
\texttt{*   **Grape:** You (4) + Me (4) = **8**} \\
\texttt{*   **Cups:** You (4) + Me (2) = 6 per cup. Two cups = **12**} \\
\texttt{*   **Hats:** You (0) + Me (1) = 1 per hat. Two hats = **2**} \\
\texttt{*   **Total:** \$8 + 12 + 2 = \textbackslash{}mathbf\{22\}\$.} \\
\\ 
\texttt{Your analysis of the **Ideal Allocation** is also spot on. To reach this maximum sum of 22, we must allocate items to the person who values them highest:} \\
\texttt{1.  **Grape:** Both value it at 4. It doesn't affect the *sum* who gets it, but it affects the fairness of the split.} \\
\texttt{2.  **Cups:** You value them at 4, I value them at 2. **You must get both cups** to maximize the sum.} \\
\texttt{3.  **Hats:** You value them at 0, I value them at 1. **I must get both hats** to maximize the sum.} \\
\\ 
\texttt{So, the only variable left is the **Grape**.} \\
\texttt{*   **Scenario A (You get Grape):** You get \{1 Grape, 2 Cups\} = \$4 + 8 = 12\$. I get \{2 Hats\} = \$2\$. Total = 14.} \\
\texttt{*   **Scenario B (I get Grape):** You get \{2 Cups\} = \$8\$. I get \{1 Grape, 2 Hats\} = \$4 + 2 = 6\$. Total = 14.} \\
\\ 
\texttt{Both scenarios yield the maximum combined score of **14**. However, Scenario A gives you 12 points and me 2 points, while Scenario B gives you 8 and me 6. Since our goal is to maximize the *sum*, either works, but Scenario A is more efficient for you personally while maintaining the max sum.} \\
\\ 
\texttt{Given that you proposed a split that resulted in a low score (10), and I have identified the optimal split that yields 14, I should now make a **secret proposal** to lock in this optimal division.} \\
\\ 
\texttt{I will propose the allocation where you get the grape and both cups, and I take the hats. This ensures we hit the maximum possible score.} \\
\\ 
\texttt{{[}Proposal: 2 hat{]}} \\
            }
        }
    }
     \\ \\

    \theutterance \stepcounter{utterance}  
    & & & \multicolumn{2}{p{0.3\linewidth}}{
        \cellcolor[rgb]{0.9,0.9,0.9}{
            \makecell[{{p{\linewidth}}}]{
                \texttt{\tiny{[GM$|$GM]}}
                \texttt{proposal} \\
            }
        }
    }
    & & \\ \\

    \theutterance \stepcounter{utterance}  
    & & & \multicolumn{2}{p{0.3\linewidth}}{
        \cellcolor[rgb]{0.9,0.9,0.9}{
            \makecell[{{p{\linewidth}}}]{
                \texttt{\tiny{[GM$|$GM]}}
                [0, 0, 2]
            }
        }
    }
    & & \\ \\

    \theutterance \stepcounter{utterance}  
    & & \multicolumn{4}{p{0.6\linewidth}}{
        \cellcolor[rgb]{0.9,0.9,0.9}{
            \makecell[{{p{\linewidth}}}]{
                \texttt{\tiny{[P1$\langle$GM]}}
                \texttt{The other player has submitted a secret proposal. It is now your turn to submit a proposal of your own.} \\
            }
        }
    }
    & \\ \\

    \theutterance \stepcounter{utterance}  
    & \multicolumn{4}{p{0.6\linewidth}}{
        \cellcolor[rgb]{0.9,0.9,0.9}{
            \makecell[{{p{\linewidth}}}]{
                \texttt{\tiny{[P1$\rangle$GM]}}
                \texttt{{[}Proposal: 1 grape, 2 cups{]}} \\
            }
        }
    }
    & & \\ \\

    \theutterance \stepcounter{utterance}  
    & & & \multicolumn{2}{p{0.3\linewidth}}{
        \cellcolor[rgb]{0.9,0.9,0.9}{
            \makecell[{{p{\linewidth}}}]{
                \texttt{\tiny{[GM$|$GM]}}
                \texttt{proposal} \\
            }
        }
    }
    & & \\ \\

    \theutterance \stepcounter{utterance}  
    & & & \multicolumn{2}{p{0.3\linewidth}}{
        \cellcolor[rgb]{0.9,0.9,0.9}{
            \makecell[{{p{\linewidth}}}]{
                \texttt{\tiny{[GM$|$GM]}}
                [1, 2, 0]
            }
        }
    }
    & & \\ \\

    \theutterance \stepcounter{utterance}  
    & & & \multicolumn{2}{p{0.3\linewidth}}{
        \cellcolor[rgb]{0.9,0.9,0.9}{
            \makecell[{{p{\linewidth}}}]{
                \texttt{\tiny{[GM$|$GM]}}
                [[1, 2, 0], [0, 0, 2]]
            }
        }
    }
    & & \\ \\

\end{supertabular}
}

\end{document}
